% ═══════════════════════════════════════════════════════════════════
% ПРИЛОЖЕНИЯ I–J: ДЕРЕВО ЗАВИСИМОСТЕЙ И КАРТА ИНТЕГРАЦИИ
% ═══════════════════════════════════════════════════════════════════

\section{Дерево зависимостей результатов}\label{app:tree}

Дерево показывает, какие результаты на каких опираются. Каждая
стрелка --- логическая зависимость, использованная в доказательстве;
висячие вершины --- первичные данные (целочисленная геометрия).

\begin{center}
\small
\begin{tabularx}{\textwidth}{@{}l >{\raggedright\arraybackslash}X@{}}
\toprule
Уровень & Результаты (опираются на нижележащие) \\
\midrule
Геометрия & кривая Ферма $C_N$; квартика Клейна; Фермат-квартика
$K3$; тор --- исходные объекты \\
Топология & род (Риман--Гурвиц); решётчатый граф; точная
последовательность; симплектический базис
$\to$ леммы~\ref{lem:genus}--\ref{lem:sympl} \\
Аналитика & базис форм; бета-интеграл; замкнутая форма периода
$\to$ леммы~\ref{lem:forms},~\ref{lem:closedform};
$\Ga$-тождества (прил.~\ref{app:gamma}) \\
Перепись & проводники; $h_d$ по Мёбиусу
$\to$ прил.~\ref{app:census} \\
Теоремы стенда & Риман как тождество
(билинейное тождество + базис);
$\Ga$-нормировка (замкнутая форма + инвентарь);
поляризация (Риман--Ходжа);
индексы (Смит + перепись)
$\to$ теоремы~\ref{th:riemann}--\ref{th:indexcomp} \\
Сертификаты & A (дивизор) $\to$ B (периоды, ядро $29$);
C--D ($\mu_4$-механика $\to$ спектры представлений);
E (конечность + моновариант);
F (Крамер + LLL);
G (спектр решёток + Коула--Селберг);
H (Раманудж + SNF + $\Ga$-объём) \\
Сводки & протокол V1--V9; лестница; мультиязычная верификация;
Lean-панель \\
\bottomrule
\end{tabularx}
\end{center}

\subsection{Критические пути}

Два критических пути определяют прочность арки. \emph{Аналитический}
путь: билинейное тождество Римана $\to$ теорема~\ref{th:riemann}
$\to$ поляризация $\to$ индексы $\to$ SNF-тип $\to$ объём $1125$.
Единственная аналитическая посылка пути --- билинейное тождество,
доказанное Стоксом на фундаментальном многоугольнике; всё прочее
--- целые числа. \emph{Арифметический} путь: перепись Мёбиуса
$\to$ суммы $=g$ $\to$ показатели $c_k$ $\to$ $\Ga$-произведение
объёма $\to$ согласование с $\sqrt{\disc}$. Оба пути завершаются
одним и тем же целым числом $1125^2$ с двух независимых сторон ---
это самый сильный внутренний контроль программы.

\subsection{Независимость путей}

Пути разделяют ровно один общий вход --- целочисленную геометрию
(кривая и её симметрии). Аналитический путь использует
дифференциальную форму и интегрирование; арифметический ---
подсчёты и конечные суммы. Ошибка в одном слое не может
воспроизвестись в другом: согласование финальных инвариантов
(дискриминант, SNF, объём) возможно только при корректности обоих
путей. Это повторяет философию сертификатов на уровне всей
программы.

\section{Карта интеграции с классическими источниками}
\label{app:integration}

Программа встраивается в классическую математику следующей картой
(расширенная версия приложения~\ref{app:repro}).

\begin{center}
\small
\begin{tabularx}{\textwidth}{@{}l l >{\raggedright\arraybackslash}X@{}}
\toprule
Источник & Слой программы & Роль и статус \\
\midrule
Риман--Гурвиц & род кривых & классика; используется как
топологический каркас \\
Формула Плюккера & род плоских кривых & классика; проверена на
всех моделях \\
Риман (билинейное) & теорема~\ref{th:riemann} & классика; выведена
Стоксом в тексте \\
Вейль (1949) & разложение якобиана Ферма & рамка; согласована с
переписью $h_d$ \\
Шиода (1979) & $\rho=20$, $H^2$ Фермат-квартики & рамка; все
следствия проверены целочисленно \\
Рибет (1976) & изогения $\Jac\sim E^3$ & рамка; целочисленные
следствия ($\Delta$, $j$) вычислены \\
Коула--Селберг (1949) & $\Ga$-слой периодов & формулы G4;
остатки $<10^{-100}$ \\
Гросс (LNM 588) & нормировка CM-периодов & H4; закрытие остатка
G \\
Гросс--Рорлих (1978) & якобианы Ферма & рамка H3; каждый период
--- $\Ga$-мономиал \\
Мёбиус / вкл.-искл. & перепись $h_d$ & классика; точная
арифметика \\
Арф (1941) & чётность спин-структур & errata E8; полный перебор \\
Смит (SNF) & тип поляризации & алгоритм; целочисленный \\
LLL (1982) & фазы в полях классов & алгоритм F4, H4 \\
Фергюсон--Бейли (PSLQ) & служебное восстановление показателей &
только как независимая сверка выведенного \\
\bottomrule
\end{tabularx}
\end{center}

Каждый источник занимает ровно отведённую ему роль: классика ---
как каркас, алгоритмы --- как инструменты, и ни один внешний
результат не подменяет выводов программы. Все точки контакта
сертифицированы независимо: согласование с Шиодой проверено
целочисленным рангом и сигнатурой, с Рибетом --- точными
$\Delta$ и $j$, с Коула--Селбергом --- остатками $10^{-100}$.

\section{Сводка обозначений Lean-панели}\label{app:leanref}

Файл \texttt{verification/lean/HodgeLaboratory.lean} содержит
следующие блоки (полные тексты --- в файле):

\begin{center}
\small
\begin{tabularx}{\textwidth}{@{}l >{\raggedright\arraybackslash}X@{}}
\toprule
Блок & Содержание \\
\midrule
\texttt{genus\_fermat} & $g(N)=(N-1)(N-2)/2$; случаи $N=15,30$;
$N=7$ (задача 1) \\
\texttt{census\_mobius} & $h_d$ по включениям-исключениям; суммы
$91$, $406$ \\
\texttt{stand\_disc} & $3^4\cdot5^6=1265625=1125^2$ \\
\texttt{snf\_product} & произведение SNF-типа $=1265625$ \\
\texttt{bch\_radicals} & радикалы $b_{Ch}(15)$, $b_{Ch}(30)$ ---
сверка выражений \\
\texttt{termination\_lcm} & $t^*=\lcm(W/g,H/g)$; случай
$48\times48\to48$ \\
\texttt{binary\_code\_check} & числа $48,212,432,114$ \\
\texttt{klein\_integers} & $105^3=343\cdot3375$;
$j=-3375=-15^3$; $\Delta=-343$ \\
\texttt{errata\_e8} & $2^6=64$; $28+36=64$ (род 3) \\
\texttt{torus\_triple} & $i^4=1$; $i^2\ne1$; $4\pi^2$-контур
(целочисленная часть) \\
\bottomrule
\end{tabularx}
\end{center}

Каждый блок проверяется ядром независимо; суммарное время
верификации --- секунды. Lean-панель не заменяет численные
протоколы, а завершает их: целые значения фиксируются машинным
выводом, и согласованные с ними численные остатки получают
окончательную интерпретацию.
